\documentclass[10pt,a4paper]{amsart}
\usepackage[margin=19mm]{geometry}
\usepackage{amsmath,amssymb,mathtools}
\usepackage[scr=rsfs,cal=euler]{mathalfa}
\usepackage{xfrac}
\usepackage[unicode,hidelinks,bookmarks=false]{hyperref}
\newcommand{\ie}{i.e.\ }
\newcommand{\cf}{cf.\ }
\newcommand{\resp}{respectively\ }
\newcommand{\Subsection}{\subsection}
\providecommand{\nobreakdash}{\nobreak\hbox{-}}
\setlength{\emergencystretch}{2em}
\multlinegap0pt
\usepackage{amssymb}
\usepackage{mathtools}
\usepackage{bm}
\usepackage[shortlabels]{enumitem}
\usepackage{algorithm}\usepackage{algpseudocode}
\usepackage{float}
\usepackage{tikz}
\newtheorem{proposition}{Proposition}
\title{Stochastic compliance–evasion dynamics in tax models:\\ a piecewise deterministic Markov process approach}
\author{Jonas Mayr, Amira Meddah, Irene Tubikanec}
\date{Source: arXiv:2605.23919v1 (CC BY 4.0)}
\begin{document}
\maketitle

\noindent\emph{Source and licence.} Stochastic compliance–evasion dynamics in tax models:\\ a piecewise deterministic Markov process approach. Version arXiv:2605.23919v1 (CC BY 4.0), distributed by arXiv under the Creative Commons Attribution 4.0 International licence, http://creativecommons.org/licenses/by/4.0/, as recorded in the arXiv licence metadata for this exact version and shown on the abstract page https://arxiv.org/abs/2605.23919v1. Full licence notice: \texttt{licenses/arxiv-2605.23919-CC-BY-4.0.txt}.\par\smallskip

\noindent\textit{Editorial scope.} Counted unit: the article's proposition prop:audit (source0.tex lines 529--531). The article's complete original proof of it is delivered with it (source0.tex lines 533--547, one \texttt{proof} environment of 15 lines). No mathematics is written, repaired or re-derived: the statement and the proof are the article's own lines, in the article's own notation, and no retained line is substituted or re-flowed.  Retained uncounted as the source-local context the proof needs: lines 499--506.  Nothing was omitted from any retained span, so the package carries no omission record.  No retained line contains a citation, so the wrapper carries no bibliography and no bibliography entry.  No retained line contains a figure environment, an image-inclusion command or drawing code, so no image is delivered and the package declares no asset dependency.  Editorial formatting only, in addition: an AMS \texttt{amsart} preamble that restates the article's own declarations and macros used by the retained lines, namely the environments \texttt{proposition} on separate counters, together with the standard packages those lines need.  The retained environments print in this document as Proposition 1 in order of appearance, whereas the article prints the same items with the numbering of its own full document.  The complete native source of the article as distributed in the arXiv e-print for this exact version is bundled unchanged as \texttt{sources/201246/source0.tex}.
\begin{equation}\label{audit_map}
  \Psi^\alpha_j(x) =
  \begin{cases}
    x^1_j + \delta \displaystyle\sum_{\beta=2}^m x^\beta_j, & \alpha = 1, \\[1.5ex]
    (1-\delta)\,x^\alpha_j, & \alpha > 1,
  \end{cases}
  \qquad j = 1, \dots, n.
\end{equation}

\begin{proposition}\label{prop:audit}
For every $\delta \in (0,1]$ and $x \in \mathcal{X}$, the audit update $\Psi$ \eqref{audit_map} satisfies $\Psi(x)\in \mathcal{X}$. 
\end{proposition}

\begin{proof}
Let $\delta \in (0,1]$ and $x \in \mathcal{X}$. By construction, it holds that 
\begin{equation*}
    \Psi^\alpha_j(x) \geq 0, \quad \text{for all } j=1,\ldots,n \ \text{and} \ \alpha=1,\ldots,m.
\end{equation*}
Moreover, for each income class $j \in \{1,\ldots,n\}$, we have
\begin{align*}
 \sum_{\alpha=1}^m \Psi_j^\alpha(x)
&= \left(x_j^1 + \delta\sum_{\beta=2}^m x_j^\beta\right) + \sum_{\alpha=2}^m (1-\delta)x_j^\alpha \\
&= x_j^1 + \delta\sum_{\beta=2}^m x_j^\beta + \sum_{\alpha=2}^m x_j^\alpha - \delta\sum_{\alpha=2}^m x_j^\alpha \\
&= x_j^1 + \sum_{\alpha=2}^m x_j^\alpha \\
&= \sum_{\alpha=1}^m x_j^\alpha.
\end{align*}
Therefore, since $x\in \mathcal{X}$, and thus $\sum_{j=1}^n\sum_{\alpha=1}^m x_j^\alpha=1$, we also have that $\sum_{j=1}^n\sum_{\alpha=1}^m \Psi_j^\alpha(x)=1$. This implies the statement.
\end{proof}

\end{document}
